% Gerado por tcc/gerar_latex.py a partir de TCC_exo1_exo2.docx.
% Edite à vontade; uma nova geração sobrescreve este arquivo.

\chapter{Fundamentação teórica}\label{sec:2}

\section{Carga do ombro no trabalho acima da cabeça}\label{sec:2.1}

Durante atividades que exigem a manutenção dos braços elevados, o sistema musculoesquelético é submetido a demandas intensas de torque articular \cite{moeller2022,perry2007}. A contração dos músculos flexores e abdutores para vencer a gravidade induz à fadiga precoce e eleva o risco de Distúrbios Osteomusculares Relacionados ao Trabalho (DORT).

A evidência epidemiológica disponível indica relação dose-dependente entre a exposição postural e o desfecho clínico. A meta-análise de \textcite{vandermolen2017} atribui à elevação de braços e mãos a duplicação do risco de distúrbios específicos do ombro, enquanto \textcite{seidler2020} quantificam acréscimo de 21\% no risco a cada mil horas acumuladas de trabalho com as mãos acima da linha dos ombros — faixa angular na qual se intensifica a compressão subacromial sobre os tendões do manguito rotador.

O torque exigido do ombro depende do ângulo de elevação (θ). Sob condições estáticas ou semiestáticas, o momento gravitacional sobre o conjunto braço–antebraço varia proporcionalmente ao seno do ângulo articular, com pico quando o segmento fica perpendicular ao vetor da gravidade, a 90° de elevação \cite{grazi2020,palazzi2022}. É esse perfil que qualquer dispositivo de apoio precisa acompanhar.

\section{Exoesqueletos de membro superior}\label{sec:2.2}

Avanços na robótica e na biomecânica aceleraram o desenvolvimento de exoesqueletos que interagem com a estrutura musculoesquelética humana \cite{gull2020}. Na reabilitação, apoiam pacientes com sequelas neurológicas, como as do Acidente Vascular Cerebral \cite{barsotti2015,bouteraa2016,zeiaee2017}; na indústria, reduzem a sobrecarga física de trabalhadores \cite{moeller2022,spada2017}.

Quanto à fonte de energia empregada, tais dispositivos classificam-se em ativos, que incorporam atuadores, sensores e malhas de controle \cite{perry2007,rahman2015}; passivos, que armazenam e restituem energia por meio de elementos elásticos, sem qualquer fonte externa de potência; e semipassivos, que regulam mecanicamente o nível de assistência sem, contudo, injetar energia no movimento \cite{grazi2020}. No contexto ocupacional de apoio ao ombro, as arquiteturas passivas predominam, e o estado da arte, os desafios abertos e as direções futuras de pesquisa desse grupo encontram-se sistematizados por \textcite{reyes2023}.

No caso dos dispositivos ativos, o desempenho é condicionado tanto pela arquitetura mecânica quanto pela instrumentação responsável por estimar a intenção de movimento do usuário. \textcite{atia2018} acompanham a cinemática do braço por unidades inerciais em protótipo de baixo custo, ao passo que \textcite{obukhov2025} recorrem à fusão de eletromiografia, unidades inerciais e rastreadores de realidade virtual para classificar a atividade motora que comanda o exoesqueleto. Tal aparato instrumental inexiste nos dispositivos passivos, nos quais a assistência fica integralmente determinada na etapa de projeto do mecanismo — circunstância que confere ao modelo preditivo um peso que ele não teria em um sistema controlado.

Os efeitos documentados pela literatura experimental apresentam direção consistente, porém magnitude variável: a meta-análise de \textcite{bar2021} reúne as reduções de esforço físico agudo associadas ao uso de exoesqueletos ocupacionais, e \textcite{moeller2022} registram diminuição da atividade de deltoide e trapézio na maioria dos 35 estudos revisados, acompanhada, entretanto, de efeitos indesejados quando o ajuste do dispositivo não corresponde à tarefa executada.

\section{Compensação de gravidade em mecanismos passivos}\label{sec:2.3}

A compensação de gravidade é problema clássico da síntese de mecanismos. \textcite{arakelian2016} sistematiza as soluções em três grupos, conforme a natureza da força compensadora: contrapesos, molas ou força ativa de um atuador auxiliar. \textcite{nguyen2026} atualiza o panorama das alternativas estritamente passivas e de seus limites de aplicação. Em exoesqueletos vestíveis, o contrapeso é inviável, e predomina a solução elástica, frequentemente associada a cabos, polias ou cames que moldam a curva de torque.

A patente US 9.889.554 B2 \cite{vanengelhoven2018} descreve um arranjo desse tipo: um elo proximal, solidário a uma base fixada ao tronco, articula-se com um elo distal que acompanha o braço, e um gerador de força de tração, formado por mola e cabo tensionado sobre uma polia, produz o torque de apoio. Abaixo de um ângulo de transição (\textit{toggle angle}), uma protuberância junto à junta restringe o gerador de força, o torque cai a um valor residual e o braço recupera a liberdade de movimento.

A patente EP 3 504 033 B1 \cite{ottobock2023} trata o mesmo problema por outro caminho construtivo: a força de um atuador passivo é transmitida por um elemento rígido apoiado na pelve, e o elemento elástico atua diretamente sobre uma alavanca solidária ao braço, sem o arranjo de cabo e polia. A faixa sem torque com o braço baixado também não depende de uma protuberância: a partir de um ângulo ajustável, o elástico passa a se apoiar num elemento de encosto junto à articulação; com esse elemento sobre o eixo de rotação, as forças dos dois trechos do elástico passam pelo eixo e não produzem torque, embora o elástico continue tracionado \cite{ottobock2023}.

Nenhum dos dois documentos, contudo, permite quantificar o torque: as curvas publicadas têm a abscissa em graus, mas a ordenada não tem unidade declarada na US 9.889.554 B2 e não tem sequer valores na EP 3 504 033 B1. Isso inviabiliza o aproveitamento direto dos dados publicados e motiva a reconstrução geométrica adotada neste trabalho.

\section{Curva torque × ângulo e ajuste ao perfil gravitacional}\label{sec:2.4}

O mapeamento da curva torque × ângulo gerada pelo mecanismo permite alinhar o perfil de suporte elástico à curva anatômica de carga do braço, otimizando a entrega de torque nas posições de maior exigência \cite{okabe2026}. \textcite{palazzi2022} mostram, em um exoesqueleto balanceado por gravidade, redução de 89\% no torque de pico do ombro quando esse ajuste é bem resolvido.

O ajuste, entretanto, encontra limite intrínseco à natureza passiva do dispositivo. \textcite{nasr2022} otimizam mecanismos passivos de ombro mediante simulação dinâmica inversa e equações algébrico-diferenciais e concluem que um mecanismo dessa classe entrega torque específico para cada ângulo de elevação, condição que o torna adequado a uma tarefa e inadequado às demais. Daí decorrem as soluções semipassivas de assistência regulável \cite{grazi2020} e as configurações ativo-passivas \cite{nasr2023}; na primeira direção, \textcite{grazi2022} regulam o nível de assistência em tempo real a partir dos sinais cinemáticos captados pelos próprios sensores embarcados no dispositivo.

A EP 3 504 033 B1 explicita esse requisito na forma de preferências de projeto: nenhuma força abaixo de um ângulo do braço em relação à vertical, como 45°, 60° ou 90°, e força máxima entre 70° e 120°, de preferência em 90° \cite{ottobock2023}. A patente destaca, nas reivindicações, dois meios de ajustar a força, a pré-tensão do elemento elástico e a excentricidade do ponto de aplicação, e publica curvas para cada um deles; a descrição cita ainda o comprimento da mola sem carga, a posição do elemento de encosto junto à articulação e a troca de atuadores.

\section{Compatibilidade cinemática e forças de interface}\label{sec:2.5}

O membro superior tem arquitetura anatômica complexa, com destaque para o complexo do ombro, modelável como junta de três graus de rotação cujo centro glenoumeral se desloca durante a elevação do braço \cite{gopura2009,shen2020}. Os desalinhamentos entre os eixos de rotação da estrutura e as articulações biológicas geram hiperestaticidade e, com ela, forças parasitas de cisalhamento e pressão excessiva nos apoios de fixação, comprometendo conforto e eficiência \cite{zeiaee2017}.

\textcite{naf2018} revisam nove categorias de estratégias de compensação de desalinhamento e as avaliam segundo o domínio de aplicação e os requisitos específicos de cada um, ao passo que \textcite{yoon2022} demonstram que a otimização da posição das juntas de um exoesqueleto de ombro reduz simultaneamente a força de interação e o volume ocupado pela estrutura. Tais forças de contato constituem, precisamente, uma das grandezas de saída que o modelo desenvolvido neste trabalho se propõe a estimar.

A EP 3 504 033 B1 trata a compatibilidade cinemática por outro caminho. A patente atribui ao estado da técnica o alinhamento dos eixos do dispositivo aos do ombro e observa que, mesmo assim, movimentos do tronco, da coluna e da cintura escapular não são acompanhados \cite{ottobock2023}. Em vez disso, o elemento de apoio do braço liga-se ao elemento de transmissão de força por uma única articulação, cujo eixo não coincide com o de abdução do ombro, e esse elemento pode girar em relação ao cinto pélvico. Com isso, a concha do braço tem três graus de liberdade de translação e três de rotação em relação ao contra-apoio, e a reação é levada à pelve, o que, segundo a patente, evita a compressão de partes moles e das costelas.

\section{Modelagem dinâmica multicorpos}\label{sec:2.6}

A análise de equilíbrio estático dimensiona a compensação de gravidade \cite{palazzi2022}, mas durante o movimento real ela é sobreposta por componentes dinâmicas: matriz de inércia M(θ), forças de Coriolis e centrífugas C(θ,θ'), atrito articular não linear com folga (\textit{backlash}) e histerese nos mecanismos de transmissão \cite{rahman2015,wang2022}. Daí a necessidade de estender o modelo a uma análise multicorpos, pela formulação de Euler-Lagrange ou por equações algébrico-diferenciais, avaliando o comportamento sob velocidades e acelerações transitórias.

Os parâmetros inerciais dos segmentos corporais que alimentam essa formulação provêm de tabelas antropométricas: \textcite{deleva1996} ajusta os valores de Zatsiorsky e Seluyanov aos centros articulares, procedimento que viabiliza o transporte de massa e momentos de inércia ao eixo da junta do exoesqueleto pelo teorema dos eixos paralelos.

O MBDyn, programa livre de dinâmica multicorpos desenvolvido no Politecnico di Milano, integra sistemas dessa natureza segundo a formulação apresentada por \textcite{masarati2014} e admite forças externas trocadas por \textit{socket}, o que viabiliza o acoplamento com um modelo implementado em Python. Trabalhos recentes percorrem esse caminho: \textcite{nasr2022} e \textcite{nasr2023} modelam o conjunto humano–exoesqueleto como cadeia cinemática fechada, enquanto \textcite{musso2022} descrevem matematicamente um exoesqueleto passivo não rígido e o avaliam por simulação musculoesquelética.

\section{Arquiteturas passivas de apoio ao ombro}\label{sec:2.7}

Três arquiteturas comerciais consolidaram soluções distintas para esse mesmo problema. O shoulderX, desenvolvido na Universidade da Califórnia em Berkeley, gera o torque por conjunto mola–cabo–polia montado entre um elo proximal e um elo distal: a mola traciona um cabo que contorna a polia solidária ao elo distal, e uma protuberância junto à junta restringe o gerador de força além do ângulo de transição, instaurando a zona neutra em que o torque é praticamente nulo com o braço baixado \cite{vanengelhoven2018}.

O Airframe, da Levitate Technologies, adota sistema adaptativo de apoio no qual o elemento elástico é progressivamente engajado à medida que o braço se eleva, transferindo-se a carga ao tronco e à cintura pélvica \cite{doyle2017}. Já o Paexo, da Ottobock, emprega elemento elástico atuando sobre braço de alavanca ajustável, de modo que o suporte é máximo com o braço próximo da horizontal e se anula com o braço junto ao corpo \cite{ottobock2023}. Apesar das diferenças construtivas, as três soluções perseguem idêntico objetivo funcional: sintetizar, por meio da geometria e de um único elemento elástico, perfil de torque que acompanhe o momento gravitacional do membro superior ao longo da elevação.

Nas configurações das Figuras 1 e 18 da patente da Ottobock, a força de apoio é transmitida por uma barra que vai da articulação junto ao ombro a um cinto pélvico, e o elemento elástico, uma mola de tração ou um cabo elástico, liga essa barra a uma alavanca do elemento que acompanha o braço, sem cabo de transmissão nem polia; a descrição admite outras variantes, como o princípio do cabo Bowden e o uso de cames ou engrenagens \cite{ottobock2023}.

A avaliação experimental desses dispositivos evidencia, contudo, que o ganho não é automático. \textcite{kim2018a} relatam redução de até 45\% no pico de atividade muscular do ombro em tarefas simuladas de furar e montar acima da cabeça, acompanhada de queda no tempo de execução e de aumento no número de erros; em trabalho complementar, os mesmos autores documentam efeitos não previstos sobre a cinemática do ombro, o equilíbrio postural e a carga na coluna \cite{kim2018b}. A leitura de projeto que daí decorre é direta: o benefício depende menos da magnitude do torque fornecido do que de sua distribuição angular, o que converte o requisito de projeto em uma curva τ(θ), e não em valor isolado de torque máximo.

\section{Balanceamento estático por molas}\label{sec:2.8}

A síntese de um perfil de torque por meio de molas dispõe de base teórica consolidada. \textcite{herder2001} formaliza os sistemas de energia nula, nos quais a soma das energias potenciais gravitacional e elástica permanece constante em toda a faixa de movimento, e demonstra que o balanceamento perfeito exige molas de comprimento livre nulo, obtidas na prática por arranjos de cabo e polia.

\textcite{wu2016} aplicam o princípio a um exoesqueleto de membro superior e relatam redução de 34,6\% no consumo de potência em relação à configuração não balanceada. Em exoesqueletos vestíveis, porém, o balanceamento perfeito não constitui o alvo: as restrições de empacotamento, o curso admissível da mola e a exigência de anular o suporte com o braço baixado convertem o problema em exercício de projeto restrito, no qual se busca a melhor aproximação de um perfil desejado dentro de limites construtivos, e não a compensação integral do peso. Por essa razão, o torque de um mecanismo passivo resulta do produto entre a força elástica, função da extensão da mola, e o braço de momento do cabo, ambos dependentes da configuração de forma não linear.

No dispositivo da EP 3 504 033 B1, o elástico liga diretamente a barra à alavanca, sem o arranjo de cabo e polia que, segundo \textcite{herder2001}, aproxima a mola da condição de comprimento livre nulo. O perfil de torque resulta da variação simultânea do comprimento do elástico e do braço de alavanca com o ângulo, e o contato com o elemento de encosto introduz uma segunda configuração geométrica, de torque nulo ou reduzido, na faixa de braço baixado.

\section{Parâmetros inerciais do membro superior e transporte de eixos}\label{sec:2.9}

O cálculo do torque gravitacional e das equações de movimento pressupõe o conhecimento da massa, da posição do centro de massa e dos momentos de inércia de braço, antebraço e mão. \Textcite{deleva1996} ajusta os dados de Zatsiorsky e Seluyanov, obtidos por varredura de raios gama, de modo que as posições de centro de massa e os raios de giração fiquem referenciados aos centros articulares usualmente adotados em biomecânica, e não aos marcos ósseos originais. Os valores são expressos como frações da massa corporal total e do comprimento do segmento, o que permite escalar o modelo para o indivíduo de referência adotado.

Como tais parâmetros são fornecidos no centro de massa de cada segmento, sua transposição ao eixo da junta do exoesqueleto demanda o teorema dos eixos paralelos, segundo o qual o momento de inércia em eixo paralelo corresponde à soma do momento no centro de massa com o produto da massa pelo quadrado da distância entre os eixos. Cabe distinguir, nesse ponto, as duas grandezas que o procedimento alimenta: o momento estático do peso, dependente apenas da massa e do braço horizontal do centro de gravidade, e o momento de inércia, que descreve a resistência à aceleração angular e comparece exclusivamente na análise dinâmica.

\section{Identificação de parâmetros a partir de curvas publicadas}\label{sec:2.10}

Quando os parâmetros construtivos não são divulgados pelo fabricante, resta identificá-los a partir daquilo que a patente efetivamente publica. As curvas de torque dessas fontes são, em geral, qualitativas: a ordenada aparece rotulada apenas como torque, sem unidade declarada, circunstância que permite ajustar a forma do perfil, mas não convertê-lo em newton-metro sem dado adicional documentado.

O procedimento usual consiste no ajuste por mínimos quadrados dos parâmetros desconhecidos contra os pontos extraídos da curva, sob restrições geométricas e de projeto explicitadas, avaliado por indicadores como a raiz do erro quadrático médio e o coeficiente de determinação. Dois cuidados distinguem um ajuste defensável de um exercício circular: a separação entre a etapa de identificação e a etapa de uso, com o congelamento dos parâmetros antes da avaliação do modelo, e a classificação de cada grandeza quanto à origem — medida, de literatura, identificada, derivada ou assumida — de modo que nenhum parâmetro ajustado seja apresentado como medida.

Quando uma mesma figura reúne curvas do mesmo dispositivo que diferem por um único parâmetro, as razões entre elas eliminam os fatores comuns. Na EP 3 504 033 B1, três curvas foram obtidas com o mesmo atuador e pré-tensões diferentes; admitindo-se que a rigidez, a escala da ordenada, o braço de momento e o torque do peso sejam comuns às três, a razão entre duas curvas depende apenas do comprimento do elástico e das pré-tensões. Essa propriedade permite identificar o comprimento de referência do elástico sem conhecer a escala da ordenada e é a base do procedimento da seção~\ref{sec:3B.9}.

\section{Lacuna}\label{sec:2.11}

A literatura revisada avalia, predominantemente, dispositivos já construídos, por meio de eletromiografia e questionários, ou otimiza mecanismos concebidos pelos próprios autores. Configura-se, assim, a ausência de uma comparação quantitativa entre os mecanismos passivos que efetivamente alcançaram o mercado, construída a partir da documentação pública de patente e conduzida do equilíbrio estático até a dinâmica multicorpos. É essa a contribuição pretendida pelo presente trabalho.
